\documentclass[12pt]{article}

\usepackage[margin=1in]{geometry}
\usepackage[T1]{fontenc}
\usepackage[utf8]{inputenc}
\usepackage{newtxtext}
\usepackage{microtype}
\usepackage{setspace}
\usepackage{xurl}
\usepackage{hyperref}
\usepackage{fancyhdr}

% ---- Document metadata ----
\newcommand{\PaperTitle}{Reading ``Problem Students'' Relationally: Deconstructing Bias Toward Neurodivergent and Culturally/Linguistically Diverse Learners Through Practicum Observation}
\newcommand{\AuthorName}{Piter Garc\'ia}
\newcommand{\Affiliation}{University of Rochester, Warner School of Education}
\newcommand{\CourseName}{EDU 442}
\newcommand{\InstructorName}{Dr. O'Neil-White}
\newcommand{\DueDate}{April 13, 2026}
\newcommand{\AssignmentName}{Deconstructing Bias \textemdash{} Reflective Paper}

\hypersetup{
  hidelinks,
  pdftitle={\PaperTitle},
  pdfauthor={\AuthorName},
  pdfsubject={\CourseName\space\textemdash{} \AssignmentName},
}

\doublespacing
\setlength{\parindent}{0.5in}
\setlength{\parskip}{0pt}
\setlength{\emergencystretch}{3em}

\pagestyle{fancy}
\fancyhf{}
\fancyhead[R]{\thepage}
\renewcommand{\headrulewidth}{0pt}
\setlength{\headheight}{15pt}

\newcommand{\levelone}[1]{\begin{center}\textbf{#1}\end{center}}
\newcommand{\reference}[1]{\noindent\hangindent=0.5in\hangafter=1 #1\par}
\newcommand{\EvidenceSet}[1]{\textsc{E#1}}
\newcommand{\JournalArtifact}[1]{\textsc{J#1}}

\begin{document}

\begin{titlepage}
\thispagestyle{fancy}
\begin{center}
\vspace*{0.14\textheight}
{\Large\bfseries \PaperTitle\par}

\vspace{0.06\textheight}
{\AuthorName\par}
{\Affiliation\par}

\vspace{1.25\baselineskip}
{\AssignmentName\par}

\vspace{1.25\baselineskip}
{\CourseName\par}
{\InstructorName\par}
{\DueDate\par}
\vfill
\end{center}
\end{titlepage}

\setcounter{page}{2}

\levelone{\PaperTitle}

This paper examines a bias that has followed me into classroom practice: the tendency to read visible dysregulation, loudness, noncompliance, or nonstandard participation as evidence of a student's fixed character rather than as evidence of context. More specifically, I examine the ``problem student'' stereotype as it attaches to learners whose difference is read through neurodivergence, culture, language, or some combination of the three. The community I observed is my elementary practicum classroom during spring 2026. The setting is in-person classroom instruction and mentor feedback conversations documented through verbatim, time-stamped observation records. Those materials are organized in this paper as Evidence Sets \EvidenceSet{1}--\EvidenceSet{4} and cross-referenced in \path{artifacts/evidence_index.md}. I also draw on Journal Artifacts \JournalArtifact{1}--\JournalArtifact{3} to make my positionality, commitments, and paper throughline transparent.

My guiding question is straightforward: \textbf{How do my practicum observations challenge the stereotype that certain neurodivergent and culturally/linguistically diverse (CLD) students are inherently disruptive or unmotivated, and what changes when adults respond relationally rather than punitively?} This question matters for future educational practice because bias does not usually appear in dramatic declarations; it appears in the small interpretive moves teachers make under pressure. A student raises their voice, moves differently, fails to comply immediately, or participates in ways that do not fit dominant discourse norms, and the adult around them has to decide what that behavior means. My concern in this paper is with that interpretive moment.

The argument I advance is that what schools often label as ``behavior'' is frequently communication: a response to unclear expectations, cognitive overload, unmet regulation needs, linguistic mismatch, public shame, or the accumulation of adult control. Across the evidence sets, I show that students positioned as problems are not naturally or permanently disruptive. Rather, the category itself is produced and reinforced through adult language, task design, and normative assumptions about what ``good'' participation is supposed to look like. When those conditions change through co-regulation, relational trust, competence-talk, and more linguistically just participation structures, different capacities come into view. In other words, the paper argues for an interpretive shift from reading behavior as defiance to reading it as evidence.

\levelone{Autobiographical Grounding: Naming the Bias}

This paper is not written from a neutral or detached standpoint. I am a neurodivergent Dominican-American student teacher who has spent much of my educational life learning what institutions reward, what they punish, and which kinds of difference they quietly translate into deficit. That history does not place me outside bias. It means, instead, that I have a particular set of instincts: I am often protective of students who are publicly misread, but under time pressure I can still feel the pull of control, efficiency, and compliance.

Talusan (2022) describes identity-conscious practice as a discipline of noticing which assumptions we carry into educational spaces and how those assumptions shape what we see in students. That framing is especially useful here because the stereotype I am examining is not a belief I consciously endorse; it is a reflexive interpretive pattern shaped by schooling itself. In many school contexts, smooth compliance is rewarded as competence, while visible difficulty is read as attitude, disrespect, or a lack of motivation. Once that association becomes habitual, adults do not merely notice behavior; they narrate a child through it.

My own identity artifacts help clarify why this paper matters to me. In my self-portrait (\JournalArtifact{1}), I described my identities as interlocking rather than fragmented: Dominican heritage, first-generation class position, ADHD, chronic illness, and culturally and linguistically diverse experience form a whole rather than a list of separate hardships. That metaphor matters because the ``problem student'' stereotype is also a story about fragmentation: it reduces a learner to the one behavior that makes adults uncomfortable. Likewise, in my journal letter to my future teacher-self (\JournalArtifact{2}), I worried about systems that claim equity while still rewarding compliance and punishing difference. I wrote there that another CLD, neurodivergent child should not be allowed to ``slip through cracks that were never accidents'' (\JournalArtifact{2}). That line remains one of the ethical anchors of this paper.

At the same time, I want to be careful not to romanticize my standpoint. Being neurodivergent and culturally/linguistically diverse does not automatically make my reading of other students accurate or just. It gives me reasons to remain suspicious of deficit narratives, but it does not relieve me of the responsibility to analyze my own habits of interpretation. The work of deconstructing bias, then, is not only about rejecting harmful labels in theory. It is about slowing down the speed with which those labels become plausible in practice.

\levelone{Observation Context and Evidence Base}

The evidence for this paper comes from four observation-based records and three identity/journal artifacts. The core observational materials are all unpublished records, so I cite them in-text as \EvidenceSet{1}--\EvidenceSet{4} rather than placing them in the References list. For transparency, each evidence set is indexed in \path{artifacts/evidence_index.md} with source file locations, summary notes, and line-level references.

The primary observation is \EvidenceSet{1}, a grade 5 media and technology lesson from March 3, 2026. During this lesson, students were creating short GIF-style videos on Chromebooks. One focal student, anonymized in the evidence index as \textit{Student M}, was repeatedly positioned through public behavioral language such as ``screaming'' and commands to put their head down (\EvidenceSet{1}, L82--L86). That same lesson, however, also includes a contrasting set of adult moves: the student is invited to help peers, receives negotiated reinforcement, and is later described by an adult with visible surprise as ``so smart'' (\EvidenceSet{1}, L72, L136, L176).

I also draw on \EvidenceSet{2}, a grade 2 robotics and routines block from February 26, 2026. This record is valuable because it captures the same student being narrated in contradictory ways within a single block of time. Early in the lesson, an adult tells the student, ``You are ruining this for your whole class'' and ``You're doing weird stuff over there'' (\EvidenceSet{2}, L4). Later, the student is repositioned through competence-talk: ``You're smart, you can help them'' (\EvidenceSet{2}, L68). Later still, the class hears an explicitly asset-based statement about cognitive diversity: ``the way that you guys think differently is a strength'' (\EvidenceSet{2}, L160).

A third source, \EvidenceSet{3}, is a mentor feedback conversation from April 8, 2026. This record matters because it provides an adult, school-based interpretation that supports a relational reading of student behavior. My mentor remarks that the patience and connection built with \textit{Student M} changed what was possible during the school day and notes that the student is ``used to people just yelling at him all day'' (\EvidenceSet{3}, L26). In the same conversation, I characterize the student's escalation as an effort to control what feels uncontrollable, something that ``triggers him'' (\EvidenceSet{3}, L32).

Finally, I use \EvidenceSet{4}, an EDU 442 seminar discussion on linguistic justice from March 9, 2026, to extend the analysis beyond behavior narrowly defined. In that seminar, participants named how schools often demand silence, speed, and standardized speech as proof of respect and learning, even though those norms do not reflect all communities equally (\EvidenceSet{4}, L236). The seminar also offered a practical counterexample: inviting students to explain coding concepts in their own language (\EvidenceSet{4}, L246). That discussion sharpened the CLD dimension of my paper by showing that behavioral bias and linguistic bias are not separate problems; both emerge when dominant norms are mistaken for neutral expectations.

Ethically, all students are anonymized, and identifying details are minimized. Names in quoted excerpts are replaced with bracketed pseudonyms. Ellipses indicate omissions for brevity, but the phrasing of quoted speech is otherwise preserved. Because this paper is based on actual practicum records rather than invented scenarios, my aim is not to diagnose students from afar, but to analyze how adults interpret and structure their participation.

\levelone{From Myth to Context: How the ``Problem Student'' Gets Produced}

The myth I am interrogating is simple, familiar, and dangerous: some students are difficult because they are naturally disruptive, disrespectful, or unwilling to learn. In school settings, that myth often rests on two assumptions. First, adults treat compliance as the clearest evidence of learning. Second, they treat deviation from dominant behavioral and discourse norms as evidence of personal deficiency. What follows is not just classroom management; it is a moral story about who belongs easily and who must be managed.

DisCrit helps explain why this story matters. Annamma, Connor, and Ferri (2013) argue that schools do not merely respond to difference; they actively produce categories of difference through sorting, labeling, and management. Read through that lens, the ``problem student'' is not simply found; the category is interactionally assembled. Flynn's (2024) affirmative non-tragedy model deepens that point by rejecting deficit and pity as the primary frames for disability or impairment. Together, these frameworks make it harder to treat regulation differences or nonnormative participation as innocent classroom inconveniences.

My practicum evidence supports this interpretation. In \EvidenceSet{1}, the focal student is publicly addressed through commands and behavioral language: ``Put your head down'' and ``you're screaming at the top of your lungs'' (\EvidenceSet{1}, L82--L86). These are not neutral observations. They situate the student's body and voice as the central problem to be corrected in public. In \EvidenceSet{2}, the pattern is even more explicit: ``You are ruining this for your whole class'' and ``You're doing weird stuff over there'' (\EvidenceSet{2}, L4). The wording matters. It does not describe a temporary difficulty; it casts the child as a threat to the collective activity.

What becomes visible, however, is that this deficit narrative is unstable. In the same grade 5 record, the adult who initially navigates disruption later recruits the student as a helper: ``I need you to help me with the kids'' (\EvidenceSet{1}, L72). A few minutes later, that adult reinforces the student's re-engagement with a negotiated reward and eventually narrates a changed perception: ``I'm actually so surprised how how smart he is'' (\EvidenceSet{1}, L136, L176). In \EvidenceSet{2}, the same learner who was framed as ``ruining'' the class is later positioned as socially useful and intellectually capable: ``You're smart, you can help them'' (\EvidenceSet{2}, L68). If the student's identity were truly fixed, such shifts should not be so immediately possible.

What the records show instead is that the ``problem student'' category is relationally produced. Public compliance demands, public shame, and a narrow interpretation of acceptable participation intensify dysregulation. Responsibility, invitation, patience, and competence-talk reopen access to the task. The student did not become a different person halfway through the lesson. The conditions around the student changed.

\levelone{Why Neurodivergence and CLD Matter to This Analysis}

I do not want the paper to collapse all forms of difference into one category. Neurodivergent learners and CLD learners are not interchangeable, and not every student in the transcripts is identified in either way. At the same time, the evidence suggests that they are often subjected to related interpretive errors. Students who move, speak, regulate, or process differently are routinely measured against narrow norms of orderly participation. Once those norms are treated as universal, difference becomes easy to misread.

The neurodivergent dimension of the paper appears most clearly in the way regulation is interpreted. In \EvidenceSet{3}, my mentor describes how patience and connection changed what was possible for \textit{Student M}: ``you've created a connection with that kid'' and ``He's used to people just yelling at him all day'' (\EvidenceSet{3}, L26). In the same conversation, I frame the student's behavior not as willful disrespect but as an effort to ``control something that he cannot control'' (\EvidenceSet{3}, L32). This does not remove responsibility from the student. It relocates the question. Instead of asking only how to stop the behavior, it asks what the behavior is doing, signaling, or protecting.

The CLD dimension becomes sharper through \EvidenceSet{4}. In that seminar, one participant described the ``demand for total silence and compliance'' as culturally loaded rather than neutral (\EvidenceSet{4}, L236). Another participant offered an example of students explaining coding in their own language (\EvidenceSet{4}, L246). These comments helped me see that schools often confuse dominant language norms with academic seriousness. A student who learns through dialogue, overlap, translation, code-switching, or nonstandard varieties can be misread as inattentive or disrespectful when, in fact, they are participating through a different communicative logic.

This is where Paris and Alim's (2014) culturally sustaining pedagogy becomes essential. If pedagogy sustains only one language variety, one discourse norm, and one mode of self-regulation, then it is not actually sustaining students; it is sustaining institutional comfort. hooks (1994) offers a related insight through engaged pedagogy: teaching that respects whole people cannot depend on reducing students to the traits that are easiest to control. The question is not whether classrooms need structure. They do. The question is whether structure will be designed to expand participation or merely to protect adult comfort.

\levelone{Assets and Funds of Knowledge}

One of the most important consequences of the deficit lens is that it obscures what students can actually do. Gonzalez, Moll, and Amanti (2005) argue that students bring funds of knowledge from home, community, language, and lived experience that schools often fail to recognize. My practicum evidence does not give a full ethnographic account of households or communities, but it does show how quickly students' capacities come into view once adults stop narrating them primarily as problems.

In \EvidenceSet{1}, the focal student is capable of helping other students once given a role that communicates trust rather than suspicion (L72). In \EvidenceSet{2}, an adult not only positions the student as helpful but explicitly names cognitive difference as a strength: ``There's never just one way to solve a problem when you're coding ... the way that you guys think differently is a strength'' (L160). That statement matters because it shifts the classroom language from accommodation-as-compensation to difference-as-resource.

The same point applies to CLD students. Linguistic diversity is not a barrier to meaning-making; it is one of the ways meaning is made. The seminar example of inviting students to explain coding in their own language (\EvidenceSet{4}, L246) is a small but powerful reminder that linguistic justice is enacted instructionally. When schools legitimate multiple ways of explaining, narrating, and reasoning, they widen what counts as academic competence.

I also want to be careful here not to replace one flattening story with another. The point is not that students are secretly extraordinary if adults would only believe in them. That is another kind of romantic simplification. The point is more modest and more important: students are often ordinarily capable in ways that deficit narratives prevent adults from seeing. Asset-based practice begins by refusing the convenience of reduction.

\levelone{Implications for Future Practice}

The evidence in this paper changes what I want to do before I assign meaning to a student's behavior. First, I want to treat escalation as information before I treat it as insubordination. This means asking contextual questions quickly and habitually: What is the student communicating? What changed in the task? What expectation is unclear? What support would make participation possible in the next two minutes? That shift from moral judgment to interpretive inquiry is not permissiveness. It is pedagogy.

Second, I want to treat relationship and co-regulation as core instructional work rather than as optional extras. \EvidenceSet{1} and \EvidenceSet{3} show that purpose, patience, connection, and competence-talk materially changed the student's access to the task and to the day. If adults know that relationship changes participation, then relationship cannot be relegated to the margins of teaching.

Third, I want to practice linguistic justice as part of classroom management. \EvidenceSet{4} makes clear that silence, speed, and standardized speech are often treated as neutral when they are not. For CLD students especially, participation may look dialogic, translanguaging, or less linear than what schools typically reward. If I want a classroom that is both rigorous and humane, then I have to design participation structures that can hold multiple communicative styles without reading difference as defiance.

Fourth, I need to audit my language. One of the clearest lessons from the transcripts is that adult wording does not merely describe events; it organizes them. Telling a child they are ``ruining this for your whole class'' (\EvidenceSet{2}, L4) is not a neutral intervention. It produces shame, positions the child as the problem, and narrows what the adult can see next. By contrast, saying ``You're smart, you can help them'' (\EvidenceSet{2}, L68) does not solve every difficulty, but it opens a different relational possibility. My goal is not to become unrealistically positive; it is to become more accurate and less harmful.

Finally, I want to preserve a commitment I articulated in \JournalArtifact{2}: rigor and care are not opposites. The danger is not structure itself. The danger is structure used primarily to purchase adult comfort through compliance. What I am trying to build toward is a classroom practice that is firm, explicit, and predictable while still leaving room for regulation differences, linguistic diversity, and human complexity.

\levelone{Conclusion}

I began this paper with a bias that is easy to carry invisibly into educational work: the reflex to treat behavioral or communicative difference as evidence of fixed student character. The practicum records and course artifacts I examined gave me reason to interrupt that reflex.

Across the evidence sets, I saw that the ``problem student'' is not a stable kind of child. It is a category that adults produce and reproduce when they confuse dominant behavioral and discourse norms with neutral expectations. In \EvidenceSet{1} and \EvidenceSet{2}, students were publicly narrated through deficit language and then, in the same records, repositioned through responsibility, competence, and trust. In \EvidenceSet{3}, relational patience was named as the condition that changed what was possible for a child who was used to being yelled at. In \EvidenceSet{4}, linguistic justice made visible how easily CLD students can also be misread when schools treat silence, speed, and standardized speech as the only legitimate signs of learning.

For me, the central lesson is both analytic and practical. Analytically, behavior is often evidence before it is defiance. Practically, that means I have to build the interpretive pause into pedagogy itself: scaffold before I shame, ask context before motive, and design participation structures broad enough to recognize more than one way of being a learner. The ambition is not to become bias-free. It is to become more accountable for the meanings I assign to students and more deliberate about the conditions I create around them.

\newpage
\levelone{References}

\reference{Annamma, S. A., Connor, D., \& Ferri, B. (2013). Dis/ability critical race studies (DisCrit): Theorizing at the intersections of race and dis/ability. \textit{Race Ethnicity and Education, 16}(1), 1--31.}

\reference{Flynn, S. (2024). Critical disability studies and the affirmative non-tragedy model: Presenting a theoretical frame for disability and child protection. \textit{Disability \& Society, 39}(2), 269--290. \url{https://doi.org/10.1080/09687599.2022.2070061}}

\reference{Gonzalez, N., Moll, L. C., \& Amanti, C. (Eds.). (2005). \textit{Funds of knowledge: Theorizing practices in households, communities, and classrooms.} Lawrence Erlbaum.}

\reference{hooks, b. (1994). \textit{Teaching to transgress: Education as the practice of freedom.} Routledge.}

\reference{Nieto, S. (2004). \textit{Affirming diversity: The sociopolitical context of multicultural education} (4th ed.). Allyn \& Bacon.}

\reference{Paris, D., \& Alim, H. S. (2014). What are we seeking to sustain through culturally sustaining pedagogy? \textit{Harvard Educational Review, 84}(1), 85--100.}

\reference{Talusan, L. (2022). \textit{The identity-conscious educator: Building habits and skills for a more inclusive school.} Solution Tree Press.}

\end{document}
